\documentclass[12pt]{article}
\usepackage[utf8]{inputenc}
\usepackage[T1]{fontenc}
\usepackage[margin=1in]{geometry}
\usepackage{setspace}
\usepackage{xcolor}
\usepackage[colorlinks=true, urlcolor=blue!60!black, linkcolor=blue!60!black]{hyperref}
\setlength{\parindent}{2em}
\setlength{\parskip}{1.25em}
\renewcommand{\baselinestretch}{1.0}
\begin{document}
\begin{center}
{\LARGE \textit{Moore v. City of East Cleveland}, 431 U.S. 494 (1977)\par}
\vspace{0.8em}
{\large Case Brief\par}
\end{center}
\vspace{1em}

\section{Memory Jogger}
In a fractured 5--4 decision, the Supreme Court reversed a grandmother's criminal conviction under an East Cleveland housing ordinance that barred her from living with two first-cousin grandsons, with a four-Justice plurality holding that substantive due process protects extended-family living arrangements as deeply rooted in the Nation's history and tradition (\href{https://supreme.justia.com/cases/federal/us/431/494/}{Justia}; \href{https://www.courtlistener.com/opinion/109670/moore-v-city-of-east-cleveland/}{CourtListener}).

\section{Detailed Case Facts}
East Cleveland, Ohio, enacted a Housing Code restricting occupancy of any dwelling unit to members of a single ``family'' (\S~1351.02) and defining ``family'' in an unusually narrow and technical way (\S~1341.08 (1966)) (\href{https://supreme.justia.com/cases/federal/us/431/494/}{Justia}). Under \S~1341.08, a ``family'' was limited to the nominal head of the household, his or her spouse, the couple's unmarried children (provided those children had no children residing with them), the parents of the head of household or spouse, and \textit{at most one} dependent child (married or unmarried) together with that child's spouse and dependent children (\href{https://supreme.justia.com/cases/federal/us/431/494/}{Justia}). Separately, \S~1351.03 regulated population density directly by prescribing a minimum habitable floor area per occupant (\href{https://supreme.justia.com/cases/federal/us/431/494/}{Justia}).

Appellant Inez Moore, a 63-year-old Black grandmother, owned a two-and-a-half-story frame house containing two dwelling units in East Cleveland (\href{https://supreme.justia.com/cases/federal/us/431/494/}{Justia}). She lived in one unit with her son, Dale Moore, Sr., and her two grandsons, Dale Moore, Jr., and John Moore, Jr. (\href{https://supreme.justia.com/cases/federal/us/431/494/}{Justia}). Because Dale, Jr., and John, Jr., were first cousins rather than brothers---John, Jr., having moved in with his grandmother after his mother died when he was less than one year old---Mrs. Moore's household did not satisfy \S~1341.08, even though it complied comfortably with the floor-area limits of \S~1351.03 (\href{https://supreme.justia.com/cases/federal/us/431/494/}{Justia}).

In January 1973, East Cleveland issued Mrs. Moore a notice of violation identifying 10-year-old John Moore, Jr., as an ``illegal occupant'' and ordering her to remove him from her home (\href{https://supreme.justia.com/cases/federal/us/431/494/}{Justia}). Although East Cleveland Building Code \S\S~1311.02--1311.03 allowed residents to apply within 10 days to the Board of Building Code Appeals for a hardship variance, Mrs. Moore did not seek a variance and refused to expel her grandson (\href{https://supreme.justia.com/cases/federal/us/431/494/}{Justia}). After sixteen months, the city filed criminal charges against her (\href{https://supreme.justia.com/cases/federal/us/431/494/}{Justia}).

\section{Procedural History}
\begin{enumerate}
\item \textbf{East Cleveland Municipal Court (May 1974):} Mrs. Moore moved to dismiss the criminal prosecution, arguing that \S~1341.08 was facially unconstitutional under the Fourteenth Amendment's Due Process and Equal Protection Clauses. The trial court overruled her motion, convicted her, and sentenced her to five days in jail and a \$25 fine (\href{https://supreme.justia.com/cases/federal/us/431/494/}{Justia}).
\item \textbf{Ohio Court of Appeals, Cuyahoga County:} Affirmed the conviction after considering her constitutional claims, relying on \textit{Village of Belle Terre v. Boraas}, 416 U.S. 1 (1974) (\href{https://supreme.justia.com/cases/federal/us/431/494/}{Justia}).
\item \textbf{Supreme Court of Ohio:} Denied review and dismissed Mrs. Moore's appeal (\href{https://supreme.justia.com/cases/federal/us/431/494/}{Justia}).
\item \textbf{U.S. Supreme Court (1976--1977):} Noted probable jurisdiction, 425 U.S. 949 (1976), heard oral argument on November 2, 1976, and reversed Mrs. Moore's conviction on May 31, 1977 (\href{https://supreme.justia.com/cases/federal/us/431/494/}{Justia}; \href{https://www.oyez.org/cases/1976/75-6289}{Oyez}).
\end{enumerate}

\section{Judicial Votes}
\textbf{Vote:} 5--4 reversing the judgment of the Ohio Court of Appeals (no single majority opinion) (\href{https://supreme.justia.com/cases/federal/us/431/494/}{Justia}; \href{https://www.oyez.org/cases/1976/75-6289}{Oyez}).
\begin{itemize}
\item \textbf{Plurality Opinion (4):} Justice Powell announced the judgment of the Court and delivered an opinion joined by Justices Brennan, Marshall, and Blackmun.
\item \textbf{Concurring Opinion (2):} Justice Brennan filed a concurring opinion joined by Justice Marshall (both also joined Justice Powell's plurality opinion).
\item \textbf{Concurrence in the Judgment (1):} Justice Stevens filed an opinion concurring in the judgment on property-rights and arbitrary-zoning grounds rather than substantive due process family privacy.
\item \textbf{Dissenting Opinions (4):} Chief Justice Burger filed a dissenting opinion (exhaustion of administrative remedies); Justice Stewart filed a dissenting opinion joined by Justice Rehnquist (merits); and Justice White filed a separate dissenting opinion (merits).
\end{itemize}

\section{Holding}
\textbf{Judgment:} Reversed (\href{https://supreme.justia.com/cases/federal/us/431/494/}{Justia}).

\textbf{Rule Announced:} A municipal housing ordinance that criminalizes cohabitation among close blood relatives---such as a grandmother and her grandsons who are first cousins---violates the Fourteenth Amendment. Under Justice Powell's four-Justice plurality, when government ``slices deeply into the family itself,'' courts abandon ordinary legislative deference and carefully examine the importance of the state interests and the extent to which the regulation serves them; substantive due process protects extended-family living arrangements because the broader institution of the family is ``deeply rooted in this Nation's history and tradition'' (\href{https://supreme.justia.com/cases/federal/us/431/494/}{Justia}). Under Justice Stevens's narrower fifth vote, an unprecedented zoning ordinance barring an owner's relatives from her home bears no substantial relation to public health, safety, morals, or general welfare under \textit{Euclid v. Ambler Realty Co.}, 272 U.S. 365 (1926), and effects an unconstitutional taking of property without due process or just compensation (\href{https://supreme.justia.com/cases/federal/us/431/494/}{Justia}).

\textbf{Holding vs. Dicta:} The binding judgment invalidates East Cleveland's criminal prohibition on extended-family cohabitation, and under \textit{Marks v. United States}, Justice Stevens's property-based concurrence supplies the narrowest formal ground supporting the judgment. However, Justice Powell's ``deeply rooted in this Nation's history and tradition'' formulation---though technically a plurality rationale rather than a five-Justice holding---became the foundational framework of modern substantive due process doctrine, later adopted by majorities in \textit{Washington v. Glucksberg}, 521 U.S. 702, 721 (1997), and \textit{Dobbs v. Jackson Women's Health Organization}, 597 U.S. 215 (2022) (\href{https://repository.law.umich.edu/cgi/viewcontent.cgi?article=4069&context=facarticles}{Mich. L. Rev.}; \href{https://harvardlawreview.org/print/vol-139/the-history-and-tradition-of-substantive-due-process-in-state-constitutions/}{Harv. L. Rev.}).

\section{Analysis of Opinions}

\subsection{Plurality Opinion of Justice Powell (Joined by Justices Brennan, Marshall, and Blackmun)}
Justice Powell first distinguished \textit{Village of Belle Terre v. Boraas}, 416 U.S. 1 (1974), which upheld a zoning ban on households of more than two unrelated persons. Whereas \textit{Belle Terre} permitted any persons related by blood, adoption, or marriage to live together, East Cleveland's ordinance ``regulat[ed] the occupancy of its housing by slicing deeply into the family itself'' (\href{https://supreme.justia.com/cases/federal/us/431/494/}{Justia}). Because the ordinance intruded directly on family living arrangements, ordinary \textit{Euclid}/\textit{Belle Terre} deference was inappropriate; instead, the Court had to scrutinize both the importance of the city's asserted interests and the fit between means and ends (\href{https://supreme.justia.com/cases/federal/us/431/494/}{Justia}).

Applying heightened scrutiny, Justice Powell found only a ``tenuous relation'' between \S~1341.08 and East Cleveland's goals of preventing overcrowding, reducing traffic congestion, and protecting school budgets. The ordinance allowed a household with a half-dozen licensed drivers or a dozen school-age children from a single son, yet prohibited Mrs. Moore from housing two grandsons who were cousins---even while \S~1351.03 already capped occupancy by floor area (\href{https://supreme.justia.com/cases/federal/us/431/494/}{Justia}).

Turning to substantive due process methodology, Justice Powell acknowledged that substantive due process ``has at times been a treacherous field'' scarred by the \textit{Lochner} era, requiring judicial restraint (\href{https://supreme.justia.com/cases/federal/us/431/494/}{Justia}). Relying on Justice Harlan's dissent in \textit{Poe v. Ullman}, 367 U.S. 497, 542--543 (1961), he argued that principled limits come not from arbitrarily stopping at the ``nuclear family,'' but from ``careful respect for the teachings of history'' and ``solid recognition of the basic values that underlie our society'' (\href{https://supreme.justia.com/cases/federal/us/431/494/}{Justia}). Because the tradition of grandparents, uncles, aunts, and cousins sharing a home---especially in times of bereavement or economic need---is deeply rooted in American history, the Constitution forbids East Cleveland from ``standardizing its children---and its adults---by forcing all to live in certain narrowly defined family patterns'' (\href{https://supreme.justia.com/cases/federal/us/431/494/}{Justia}). In footnote 5, the plurality also rejected Chief Justice Burger's exhaustion objection, explaining that a criminal defendant raising a facial constitutional challenge need not first seek a discretionary zoning variance (\href{https://supreme.justia.com/cases/federal/us/431/494/}{Justia}).

\subsection{Concurring Opinion of Justice Brennan (Joined by Justice Marshall)}
Joining the plurality in full, Justice Brennan wrote separately to expose the ``cultural myopia'' of confining constitutional protection to the nuclear family model prevalent in ``white suburbia'' (\href{https://supreme.justia.com/cases/federal/us/431/494/}{Justia}). Citing sociological and census data showing that extended-family households are a vital source of economic survival and childrearing support for working-class, immigrant, and especially Black families---noting that 48\% of Black households headed by an elderly woman included related minor children other than her own offspring, compared with 10\% of white counterparts---he insisted that the Constitution does not allow government to impose suburban nuclear-family preferences on the rest of society (\href{https://supreme.justia.com/cases/federal/us/431/494/}{Justia}). He also rejected reliance on East Cleveland's variance procedure, noting that \textit{Euclid} does not require exhaustion for facial challenges and that a discretionary ``escape hatch'' cannot rescue an unconstitutional burden on fundamental rights (\href{https://supreme.justia.com/cases/federal/us/431/494/}{Justia}).

\subsection{Opinion of Justice Stevens Concurring in the Judgment}
Declining to join the plurality's substantive due process family-privacy analysis, Justice Stevens framed the case around a homeowner's traditional property right ``to use her own property as she sees fit'' (\href{https://supreme.justia.com/cases/federal/us/431/494/}{Justia}). Surveying fifty years of state zoning decisions applying \textit{Euclid v. Ambler Realty Co.}, 272 U.S. 365 (1926), and \textit{Nectow v. Cambridge}, 277 U.S. 183 (1928), he found no precedent sustaining an ordinance that barred an owner's own blood relatives from permanent residence (\href{https://supreme.justia.com/cases/federal/us/431/494/}{Justia}). Because East Cleveland completely failed to explain why two grandchildren could live with their grandmother if they were brothers but not if they were cousins, the ordinance lacked any ``substantial relation to the public health, safety, morals, or general welfare'' and constituted a taking of property without due process and without just compensation (\href{https://supreme.justia.com/cases/federal/us/431/494/}{Justia}).

\subsection{Dissenting Opinion of Chief Justice Burger}
Chief Justice Burger did not reach the constitutional merits. Invoking administrative exhaustion, constitutional avoidance (\textit{Ashwander v. TVA}), federalism, and federal court docket overload, he argued that Mrs. Moore's ``deliberate refusal'' to seek a hardship variance from East Cleveland's Board of Building Code Appeals within the 10-day window should bar her from pressing federal constitutional objections (\href{https://supreme.justia.com/cases/federal/us/431/494/}{Justia}).

\subsection{Dissenting Opinion of Justice Stewart (Joined by Justice Rehnquist)}
Justice Stewart argued that \textit{Belle Terre} largely controlled. Freedom of association under the First Amendment protects expressive and ideological groups, not residential convenience, and prior family-privacy cases (\textit{Meyer}, \textit{Pierce}, \textit{Griswold}, \textit{Loving}, \textit{Roe}) protected decisions to marry, bear, and raise one's own children---not a grandmother's desire to share ``a single kitchen and a suite of contiguous rooms'' with extended relatives (\href{https://supreme.justia.com/cases/federal/us/431/494/}{Justia}). Because no fundamental right was implicated, the ordinance triggered only rational-basis review under the Equal Protection Clause, which it satisfied as a permissible legislative line-drawing choice, reinforced by the availability of a hardship variance (\href{https://supreme.justia.com/cases/federal/us/431/494/}{Justia}). In footnote 7, Justice Stewart also pointed out that East Cleveland was a predominantly Black city with a Black City Manager and City Commission, undercutting Justice Brennan's ``white suburbia'' framing (\href{https://supreme.justia.com/cases/federal/us/431/494/}{Justia}).

\subsection{Dissenting Opinion of Justice White}
Justice White began with a sweeping institutional warning against expanding substantive due process, noting that its substantive content is suggested neither by constitutional text nor preconstitutional history and that the judiciary comes ``nearest to illegitimacy'' when creating judge-made constitutional rights (\href{https://supreme.justia.com/cases/federal/us/431/494/}{Justia}). While accepting that Mrs. Moore possessed a ``liberty'' interest under \textit{Meyer}, he contended that residing with more than one set of grandchildren is not ``implicit in the concept of ordered liberty'' under \textit{Palko v. Connecticut}, 302 U.S. 319 (1937), and criticized Justice Powell's ``history and tradition'' test as an overly broad charter for judicial invalidation of social legislation (\href{https://supreme.justia.com/cases/federal/us/431/494/}{Justia}). Applying rational-basis review under both Due Process and Equal Protection, he concluded that the ordinance rationally advanced legitimate zoning objectives (\href{https://supreme.justia.com/cases/federal/us/431/494/}{Justia}).

\section{Examples: Future Applications}
\begin{enumerate}
\item \textbf{Same-Side (Household/Family Prevails):} A municipality enacts a single-family zoning ordinance defining ``family'' solely as parents and their biological or adopted children, and prosecutes an adult woman for allowing her elderly aunt and widowed sister to live with her in a four-bedroom home. Under \textit{Moore}, the ordinance slices directly into close blood kinship without advancing density goals better than neutral square-footage limits and is unconstitutional.
\item \textbf{Same-Side (Household/Family Prevails):} A city housing code permits a homeowner to live with up to four biological grandchildren from one married child, but prohibits living with two grandchildren born to two different unmarried children. Like East Cleveland's \S~1341.08, this classification arbitrarily discriminates among close extended-family configurations and bears only a tenuous relation to overcrowding or traffic.
\item \textbf{Opposite-Side (Municipality Prevails):} A college town enforces a zoning ordinance barring more than three unrelated individuals from occupying a single-family home, while permitting any number of persons related by blood, marriage, or adoption to live together, and cites five unrelated university roommates. Because the ordinance regulates only unrelated occupants and leaves nuclear and extended families untouched, \textit{Village of Belle Terre v. Boraas}---which \textit{Moore} expressly preserved---controls, and the ordinance survives rational-basis review.
\item \textbf{Opposite-Side (Municipality Prevails):} A city enforces a content-neutral, family-blind housing-safety code requiring at least 150 square feet of habitable floor area per occupant and cites a twelve-person extended family occupying a 600-square-foot studio apartment. Both Justice Powell's plurality (footnote 7) and Justice Stevens's concurrence (footnote 16) explicitly endorsed neutral floor-area density limits as legitimate, closely tailored means of preventing overcrowding.
\item \textbf{Fence-Sitter (Close Call):} A municipality restricts single-family zones to persons related by blood, marriage, or legal adoption, and cites a long-term ``functional family'' consisting of an unmarried couple, one partner's adult sibling, and a godchild whom the couple has raised for ten years without formal adoption. Because \textit{Moore} grounded heightened protection in traditional biological and legal kinship while preserving \textit{Belle Terre}'s distinction for unrelated persons, federal courts split sharply on whether non-biological ``chosen'' or functional families receive \textit{Moore}'s protection or fall under \textit{Belle Terre}'s deferential rational-basis rule (\href{https://fordhamlawreview.org/wp-content/uploads/2017/04/Carbone-Cahn_May.pdf}{Fordham L. Rev.}).
\end{enumerate}

\section{Critique}
\textit{Moore v. City of East Cleveland} occupies a pivotal and contested place in substantive due process jurisprudence. From an \textbf{originalist and textualist perspective}, critics echo Justice White's warning that the Due Process Clause guarantees fair procedure rather than an open-ended judicial license to strike down local land-use ordinances, arguing that Justice Powell's ``deeply rooted in this Nation's history and tradition'' standard still leaves judges enormous discretion to choose the level of generality at which a tradition is defined (\href{https://supreme.justia.com/cases/federal/us/431/494/}{Justia}; \href{https://repository.law.umich.edu/cgi/viewcontent.cgi?article=4069&context=facarticles}{Mich. L. Rev.}). At the same time, later conservative originalists embraced Powell's history-and-tradition formulation in \textit{Glucksberg} and \textit{Dobbs} as the primary \textit{limiting} device to constrain unenumerated rights (\href{https://texaslawreview.org/in-defense-of-substantive-due-process/}{Tex. L. Rev.}). Some property-rights scholars also praise Justice Stevens's concurrence as a more textually grounded alternative that avoids inventing privacy rights while checking arbitrary municipal interference with residential property (\href{https://scholarship.law.gwu.edu/cgi/viewcontent.cgi?article=2515&context=faculty_publications}{GW Law Scholarly Commons}).

From a \textbf{progressive and doctrinal perspective}, scholars celebrate \textit{Moore} for rejecting an exclusionary nuclear-family baseline and recognizing the constitutional dignity of multigenerational and kinship-care households, particularly in working-class and Black communities (\href{https://scholarship.law.gwu.edu/cgi/viewcontent.cgi?article=2515&context=faculty_publications}{GW Law Scholarly Commons}). Yet progressive family-law scholars also critique \textit{Moore}'s ``Trojan-Horse quality'': by tying constitutional protection to traditional biological and marital kinship while leaving \textit{Belle Terre} intact, Powell's opinion entrenched a normative, blood-based definition of ``family'' that courts later used to deny protection to unmarried partners, foster families, and non-traditional ``chosen'' households (\href{https://fordhamlawreview.org/wp-content/uploads/2017/04/Carbone-Cahn_May.pdf}{Fordham L. Rev.}). Moreover, Justice Powell's history-and-tradition test ultimately became the doctrinal engine for curtailing newer autonomy rights in \textit{Bowers}, \textit{Glucksberg}, and \textit{Dobbs} (\href{https://texaslawreview.org/in-defense-of-substantive-due-process/}{Tex. L. Rev.}).

Underlying \textit{Moore} is a recurring constitutional tension among competing values: family integrity and personal autonomy in domestic living arrangements, local democratic self-governance and police-power zoning authority, private property rights, and judicial restraint in enforcing unenumerated liberties.

\section{Key Quotations}
\begin{quote}
``East Cleveland, in contrast, has chosen to regulate the occupancy of its housing by slicing deeply into the family itself.''
\end{quote}
\noindent --- Justice Powell (plurality opinion), \href{https://supreme.justia.com/cases/federal/us/431/494/}{\textit{Moore v. City of East Cleveland}, 431 U.S. 494, 498 (1977) (Justia)}.

\begin{quote}
``Appropriate limits on substantive due process come not from drawing arbitrary lines, but rather from careful `respect for the teachings of history [and] solid recognition of the basic values that underlie our society.' ''
\end{quote}
\noindent --- Justice Powell (plurality opinion) (quoting \textit{Griswold v. Connecticut}, 381 U.S. 479, 501 (1965) (Harlan, J., concurring)), \href{https://supreme.justia.com/cases/federal/us/431/494/}{\textit{Moore v. City of East Cleveland}, 431 U.S. 494, 503 (1977) (Justia)}.

\begin{quote}
``Our decisions establish that the Constitution protects the sanctity of the family precisely because the institution of the family is deeply rooted in this Nation's history and tradition.''
\end{quote}
\noindent --- Justice Powell (plurality opinion), \href{https://supreme.justia.com/cases/federal/us/431/494/}{\textit{Moore v. City of East Cleveland}, 431 U.S. 494, 503 (1977) (Justia)}.

\begin{quote}
``The Constitution cannot be interpreted, however, to tolerate the imposition by government upon the rest of us of white suburbia's preference in patterns of family living.''
\end{quote}
\noindent --- Justice Brennan, concurring, \href{https://supreme.justia.com/cases/federal/us/431/494/}{\textit{Moore v. City of East Cleveland}, 431 U.S. 494, 508 (1977) (Justia)}.

\begin{quote}
``The Judiciary, including this Court, is the most vulnerable and comes nearest to illegitimacy when it deals with judge-made constitutional law having little or no cognizable roots in the language or even the design of the Constitution.''
\end{quote}
\noindent --- Justice White, dissenting, \href{https://supreme.justia.com/cases/federal/us/431/494/}{\textit{Moore v. City of East Cleveland}, 431 U.S. 494, 544 (1977) (Justia)}.

\end{document}
